{\color{gray}\hrule}
\begin{center}
\section{Hypotheses and Methodology}
% TODO One sentence: what is tested, against what, and how it is validated.
\placeholder[1]{8}
\bigskip
\end{center}
{\color{gray}\hrule}
\begin{multicols}{2}
\subsection{Hypotheses}

% TODO H0 and H1 reworded: the fair test is against Allan and Williams and the count baselines, rebuilt on every fold; DeepBach is a declared reference that has seen the melodies. Directional or not, and why; operational definitions of every measure and threshold; the date on which the hypotheses were fixed.
\placeholder{9}

\subsection{Analytical Approach}

% TODO The walk, its parts and the ablation; the measures, each with its reference; parameters and how they were chosen; software and versions.
\placeholder{10}

\subsection{Validation Approach}

Every number on how well a model imitates a style comes from cross-validation by group. We split each corpus into ten folds, as Kohavi recommends \cite{kohavi1995}, and test each fold with a model counted on the other nine. A group never crosses a fold, so that the model never meets in training a piece closely related to the one it is tested on. In Bach, a group is a tune, identified by the Zahn number \cite{zahn1889} that When in Rome gives each chorale. Bach set some tunes many times, and a setting of the same tune in training would let the model recall the tune instead of imitating the style. Two chorales also join one group when their top voices contain each other's intervals at 90\% or more, which covers the chorales without an analysis, and so do the versions of one piece. Bach thus has 181 groups, the largest with 13 files on one tune. In Palestrina a group is a mass, so that the movements of one mass never fall on both sides of a split. The corpus holds 101 masses.

We assign the groups to the folds from the largest down, each to the fold with the fewest notes so far. Every fold then holds between 7,784 and 7,909 weighted notes in Bach, and between 67,270 and 70,043 in Palestrina. The assignment has no random part, so every run gives the same folds.

% TODO The null model, what it preserves and what it randomizes; the test, its statistic, more than 30 null instances, alpha; multiple testing; effect sizes; robustness checks decided in advance, among them the repetition of the main result on five shuffled group splits with fixed seeds.
\placeholder{11}
\end{multicols}
